\section{Method}
\label{sec:method}

\subsection{Overview}
\label{sec:setting}

We consider $\Nc$ clients that hold private labelled data over a common label space of $\Cc$ classes and
want a shared classifier without pooling their data. Client $j$ holds $\Dj$ of size $\nj$; the data are
heterogeneous, with each client's class proportions drawn from a Dirichlet distribution of concentration
$\alpha$ (smaller $\alpha$ means more skew). Two constraints separate our setting from ordinary federated
learning. The protocol is \emph{one-shot}: one upload and one download per client, with no iterative
exchange. And privacy is \emph{cryptographic}: a client's contribution is never revealed in the clear, and
we are unwilling to accept the accuracy loss that differential privacy incurs when it is the
only line of defence.

These two constraints, together with the untrusted setting, do not merely shape the design; they determine
it. The protocol is not an assembly of independent techniques but the forced solution to a small set of
constraints that secure collaborative training imposes, and we arrived at it by discarding alternatives that
violate one or another --- a sequential or split-training scheme, for example, hands each party the
intermediate state its predecessors trained, exactly the exposure \textbf{C5} below forbids. The
constraints, and the choice each forces, are:

\begin{itemize}
\item[\textbf{C1}] \emph{Sophisticated learning cannot run under encryption.} Nonlinear optimisation ---
  gradient steps, softmax, adaptive curvature --- needs deep, repeatedly bootstrapped homomorphic circuits.
  \emph{So} the learning happens in plaintext on the clients, which pre-compute everything the aggregation
  will consume: the local fine-tune, the per-coordinate importance estimate, and the single displacement each
  client chooses to upload. The server performs no learning.
\item[\textbf{C2}] \emph{The server cannot be adaptive.} Blind to the values it holds, it cannot inspect,
  compare, or branch, and so cannot pick a data-dependent aggregation rule. \emph{So} it emits several
  candidate aggregates and defers the choice to the clients, who select among them after decryption
  (\cref{sec:candidates}).
\item[\textbf{C3}] \emph{The parties must share a frame before they begin.} In the open-ended landscape of
  deep-network parameters, independently trained models occupy incomparable regions and their linear
  combination is meaningless. \emph{So} the parties fix a common geometry in advance: a frozen public
  pretrained backbone fixes the representation, and a common public initialisation places every client at one
  point (\cref{sec:basin}).
\item[\textbf{C4}] \emph{The server's operation must be shallow.} Homomorphic cost is dominated by
  multiplicative depth. \emph{So} the one server operation is a depth-one linear combine, and the design is
  arranged so this cheapest operation is \emph{exactly} the one the method needs --- freezing the adapter's
  down-projection makes the linear sum of the trained factors equal the sum of the weight updates
  (\cref{sec:aggregation}).
\item[\textbf{C5}] \emph{Intermediate training signals are the strongest attack surface.} Each exposed round
  is a fresh training-time artifact, far more revealing than the released model (\cref{sec:related-leakage}).
  \emph{So} the protocol is one-shot: a single encrypted exchange, one artifact, never in the clear.
\item[\textbf{C6}] \emph{No single party may be trusted to decrypt.} The parties are mutually distrustful and
  the server is untrusted; a single decryption key would let its holder reveal every contribution. \emph{So}
  the key is generated and held distributively --- multiparty CKKS with threshold decryption, under which no
  party, the server included, can decrypt alone (\cref{sec:mhe}).
\end{itemize}

\noindent The six form a chain rather than a list: \textbf{C1} pushes the work onto the clients, which raises
\textbf{C2} (a blind server) and \textbf{C3} (their returns must be comparable); \textbf{C4} and \textbf{C6}
shape the single operation the server may perform; \textbf{C5} fixes how many times it happens. The protocol
below is one instance of their joint solution, and each constraint can be relaxed only at a stated price:
abandon one-shot encryption and the intermediate snapshots leak (\textbf{C5}); abandon encryption altogether
and the guarantee weakens from cryptographic to statistical --- the differentially private route whose
accuracy cost \cref{sec:experiments} measures. Fine-tuning a pretrained model is thus not a convenience but
the learning setting these constraints admit, and it removes the alignment phase that most distinguishes
HE-OFT from prior one-shot federated learning, which spends an extra communication round, and often a
differential-privacy budget, to bring clients into a common frame before they begin.

Concretely, every client uses a frozen public backbone $\featx$ followed by a small trainable unit:
low-rank adapters together with a classifier head, with parameters $\headp$. A low-rank adapter writes the
update to a weight matrix as a product $\Delta W = BA$ of two small factors; we freeze the down-projection
$A$ at its shared public initialisation and train only the up-projection $B$ and the head. This single
choice is what makes the encrypted aggregation \emph{exact}: with $A$ fixed and common to all clients, a
weighted sum of the trained factors reconstructs the weighted sum of the updates,
$\big(\sum_j \wj B_j\big)A = \sum_j \wj B_jA$, whereas training both factors would make the update bilinear
and the per-factor average a different, and worse, quantity. The backbone never changes, is
never transmitted, and is never encrypted; only the trainable unit is fine-tuned, communicated, and operated
on under encryption. Keeping that unit small is what makes the encrypted step inexpensive: it is a tiny fraction of the backbone
(a low-rank adapter and head, on the order of $10^5$ parameters against a backbone of $10^8$), so a client's
entire encrypted contribution is a few tens of ciphertexts rather than the model itself. Every client
initializes the trainable unit to the \emph{same} public constant $\thz$, the standard adapter
initialization, agreed in advance and carrying no client data. There is no prototype exchange and no
peer-to-peer phase; nothing data-derived ever leaves a client before encryption.

\begin{figure*}[t]
\centering
\resizebox{\textwidth}{!}{%
\begin{tikzpicture}[
  font=\small,
  card/.style={draw=sanzotan!80!black, fill=sanzotanlt, rounded corners, align=center, inner sep=3pt},
  teacher/.style={draw=sanzotan!80!black, fill=sanzotan, rounded corners, minimum width=11mm, minimum height=6mm, font=\scriptsize, align=center},
  data/.style={cylinder, shape border rotate=90, draw=sanzotan!80!black, fill=sanzotangrey, aspect=0.28,
               minimum width=10mm, minimum height=8mm, font=\scriptsize, align=center, inner sep=1pt},
  basin/.style={draw=paperblue, fill=blueB, rounded corners, minimum width=22mm, minimum height=12mm, align=center, font=\small},
  server/.style={draw=sanzobluegrey!80!black, fill=sanzobluegrey, rounded corners, minimum width=30mm, minimum height=14mm, align=center, font=\small},
  glob/.style={draw=paperblue, fill=blueD, text=white, rounded corners, minimum width=20mm, minimum height=12mm, align=center, font=\small},
  arr/.style={-{Stealth[length=2.4mm]}, semithick, draw=paperneutral},
  p2p/.style={{Stealth[length=2mm]}-{Stealth[length=2mm]}, semithick, draw=paperblue},
  lbl/.style={font=\scriptsize\itshape, text=paperneutral, align=center, inner sep=1.5pt}
]
% ---- clients (three, stacked) ----
\foreach \i/\y/\name/\s in {1/2.5/{Client 1}/1, 2/0/{Client 2}/2, 3/-2.5/{Client $N$}/N} {
  \node[data]    (d\i) at (0,\y) {data\\$\mathcal{D}_{\s}$};
  \node[teacher] (t\i) at (2.0,\y) {teacher\\$t_{\s}$};
  \draw[arr] (d\i) -- (t\i);
  \node[lbl] at (1.0,\y+0.85) {\name};
}
\begin{scope}[on background layer]
  \foreach \i in {1,2,3} \node[card, fit=(d\i)(t\i)] {};
\end{scope}
% P2P exchange among clients
\draw[p2p] (t1.east) .. controls (3.4,2.5) and (3.4,0) .. (t2.east);
\draw[p2p] (t2.east) .. controls (3.4,0) and (3.4,-2.5) .. (t3.east);
\node[lbl, text=paperblue] at (3.55,1.25) {P2P\\prototypes};

% ---- shared loss basin ----
\node[basin] (basin) at (6.2,0) {shared loss\\basin $\theta_0$};
\foreach \i in {1,2,3} \draw[arr, draw=paperblue] (t\i.east) -- (basin.west);

% ---- local distillation -> encrypted displacement ----
\foreach \i/\y/\s in {1/2.5/1, 2/0/2, 3/-2.5/N} {
  \node[teacher, fill=blueC, draw=paperblue, text=white] (delta\i) at (9.2,\y) {$\mathsf{Enc}(\Delta_{\s})$};
  \draw[arr, draw=paperblue] (basin) -- node[lbl, sloped, above]{} (delta\i.west);
}
\node[lbl] at (9.2,3.5) {distil from $\theta_0$\\($K$ steps), encrypt};

% ---- server (encrypted aggregation) ----
\node[server] (srv) at (12.8,0) {Server\\$\theta_0+\sum_j w_j\Delta_j$\\\scriptsize(CKKS, depth $1$)};
\foreach \i in {1,2,3} \draw[arr] (delta\i.east) -- (srv.west);
\begin{scope}[on background layer]
  \node[draw=paperblue, dashed, rounded corners, fit=(srv), inner sep=3mm,
        label={[lbl, text=paperblue]above:encrypted --- no party decrypts alone}] {};
\end{scope}

% ---- threshold decrypt -> global model ----
\node[glob] (glob) at (16.4,0) {global\\model $\theta^\star$};
\draw[arr] (srv.east) -- node[lbl, above]{threshold\\decrypt} (glob.west);
\end{tikzpicture}%
}
\caption{HE-IFD. Each client holds private data $\mathcal{D}_j$ and trains a local teacher $t_j$. In
Phase~0 the clients exchange per-class feature prototypes over peer-to-peer channels, with the server
excluded, to agree on a shared loss basin $\theta_0$. Each client then distils its teacher into the shared
head for a bounded number of steps and encrypts the resulting displacement $\Delta_j$. The server's sole
operation is the sample-weighted linear combination $\theta^\star=\theta_0+\sum_j w_j\Delta_j$ under
multiparty CKKS (additions and public-scalar multiplications only; multiplicative depth one); it never
decrypts. A threshold of clients jointly decrypts the global model $\theta^\star$.}
\label{fig:protocol}
\end{figure*}


The protocol is two steps, shown in \cref{fig:protocol}. Starting from the shared initialization $\thz$,
each client fine-tunes its trainable unit on its own data for a bounded number of steps and records the
displacement it travelled, $\disp = \theta_j^{(\Ksteps)} - \thz$. The server then combines the encrypted
displacements into a single global model, $\thstar = \thz + \sum_j \wj \disp$, using only homomorphic
additions and multiplications by public scalars, after which a threshold of clients jointly decrypts the
result.

The protocol consists of four steps, where $\mathsf{Enc}(\cdot)$ denotes encryption under the clients'
collective CKKS key. The learning is performed entirely on the clients, and the server's only operation is
the single encrypted sum of the third step.

\emph{Setup.} The $\Nc$ clients run a distributed key-generation protocol that yields a collective CKKS
public key, under which any party may encrypt and compute but no party can decrypt alone, decryption
requiring a threshold of clients (\cref{sec:mhe}).

\emph{Local fine-tuning, performed in parallel by each client $j$.} Beginning from $\theta_j^{(0)} = \thz$,
client $j$ takes $\Ksteps$ steps of supervised fine-tuning on its own data,
\[
\theta_j^{(k)} = \theta_j^{(k-1)} - \eta\,\nabla\mathcal{L}\big(\model;\,\Dj\big), \qquad k = 1,\dots,\Ksteps,
\]
and uploads the encryption $\mathsf{Enc}(\disp)$ of its displacement $\disp = \theta_j^{(\Ksteps)} - \thz$.

\emph{Encrypted aggregation, performed by the server.} The server forms
\[
\mathsf{Enc}(\thstar) = \thz + \sum_{j=1}^{\Nc} \wj\,\mathsf{Enc}(\disp),
\]
using only multiplications of a ciphertext by the public scalar $\wj$ and additions of ciphertexts, with the
public constant $\thz$ added in the clear; the circuit has multiplicative depth one. The server broadcasts
$\mathsf{Enc}(\thstar)$.

\emph{Threshold decryption.} A threshold of clients jointly decrypts $\mathsf{Enc}(\thstar)$ by collective
key-switching to recover the global model $\thstar$. The server never obtains a plaintext.

The objective is not a model that surpasses each client on that client's own dominant classes; a local
specialist on a narrow slice of the label space is difficult to surpass on that slice, and we do not claim to. The
objective is a common model every party can use, that classifies well across the whole label space including
classes a given client never observed locally, and that is produced without any party exposing its data or
its update in the clear.

\subsection{Notation}
\label{sec:notation}

\Cref{tab:notation} collects the symbols used throughout. Client indices are $j\in\{1,\dots,\Nc\}$.

\begin{table}[t]
\centering
\caption{Notation.}
\label{tab:notation}
\begin{tabular}{ll}
\toprule
Symbol & Meaning \\
\midrule
$\Nc$ & number of participating clients \\
$\Dj,\ \nj$ & client $j$'s local dataset and its size $\nj=|\Dj|$ \\
$\Cc$ & number of classes \\
$\alpha$ & Dirichlet concentration (label heterogeneity) \\
$\featx$ & frozen public backbone (never trained or encrypted) \\
$\headp$ & parameters of the trainable unit (adapter $+$ head) \\
$\model$ & model on the frozen features, parameters $\theta$ \\
$\thz$ & shared public initialization of the trainable unit \\
$\Ksteps$ & local fine-tuning steps \\
$\disp$ & displacement $\theta_j^{(\Ksteps)}-\thz$ \\
$\wj$ & sample weight $\nj/\sum_i n_i$ \\
$\thstar$ & aggregated global model $\thz+\sum_j\wj\disp$ \\
\bottomrule
\end{tabular}
\end{table}

\subsection{Multiparty Homomorphic Encryption}
\label{sec:mhe}

We build on the CKKS scheme~\cite{cheon2017ckks}, a ring-learning-with-errors
construction~\cite{lyubashevsky2010ideal} that performs approximate arithmetic on vectors of real numbers
under encryption. CKKS packs many real slots into one ciphertext and supports two homomorphic operations of
interest here: addition of two ciphertexts, and multiplication of a ciphertext by a plaintext. It is a
levelled scheme, meaning each multiplication consumes one level of a finite budget; our protocol uses only
additions and multiplications by public plaintext scalars, so it consumes a single level and never requires
the expensive bootstrapping step~\cite{cheon2018bootstrapping} that refreshes a depleted budget.

A single-key scheme would require one party to hold the secret key and therefore be trusted with decryption.
We instead use the multiparty variant of CKKS~\cite{mouchet2021multiparty}, in which the secret key is
shared among the clients through a distributed key-generation protocol that produces a collective public key
without ever materialising the secret in one place. Any party, including the server, may compute on
ciphertexts under the collective public key, but decryption requires the cooperation of a threshold of
clients, carried out as a collective key-switching protocol. No single party, and in particular not the
server, can decrypt on its own. This is exactly the property our aggregation needs: the server performs the
homomorphic computation on encrypted contributions, yet only the clients, acting together, can reveal the
result.

We stress what the scheme is chosen for, since our circuit uses almost none of its homomorphic power and a
reader may ask why a fully homomorphic scheme is deployed at depth one. The answer is the key structure
and the data layout, not the depth. Threshold decryption over a collectively generated key is the property
the protocol cannot do without, and the mature distributed key generation and collective key-switching of
multiparty CKKS provide it directly~\cite{mouchet2021multiparty}. CKKS additionally packs thousands of
real-valued parameters into one ciphertext and multiplies them by public scalars natively, the operations
the candidate aggregates of \cref{sec:candidates} require; additively homomorphic schemes offer neither
packed real arithmetic at this width nor comparable multiparty tooling, and masking-based secure
aggregation reveals the plaintext sum to the server, which our protocol never does. Depth one is therefore
not an underuse of the scheme but the point of the design: the protocol is engineered so that the
strongest available tool is needed only at its least expensive setting.

\subsection{The Shared Frame}
\label{sec:basin}

The natural one-shot recipe (let each client train independently and have the server average the
results) fails under heterogeneity. When clients begin from different initializations and fit different
skewed distributions, their trained parameters occupy different regions of parameter space and their average
lies in none of them: the updates point in conflicting directions and cancel when summed. The remedy is for
every client to begin from the same point and move only a little, so their displacements live in one frame
and add constructively rather than fighting one another.

A pretrained backbone provides this directly. Because the representation is fixed and shared, the
trainable unit lives in the same parameter space for every client, and initializing it to a common public
constant $\thz$ places all clients at the same point in that space before any training. No data is consulted
to build $\thz$ and nothing is exchanged to agree on it; it is a fixed, public initializer. A from-scratch
federation cannot take this step. It must first agree on a starting point, in a data-dependent alignment
phase that costs an extra round and exposes a per-client summary; protecting that summary spends a
differential-privacy budget. Fine-tuning a pretrained model removes the phase, and with it the only channel
through which client data would otherwise leave in the clear before encryption. We verify in
\cref{sec:experiments} that this shared start is what makes the one-shot combination work (averaging
independently trained units without it collapses), and that the frozen backbone supplies the start with no
alignment phase.

The frame the backbone supplies is the representation, not the classifier. The shared initializer $\thz$ is
an untrained head: it places every client at the same point in parameter space but encodes nothing about any
class. A client therefore contributes a useful displacement for a class only when it holds examples of that
class, so the global model's accuracy on a class is bounded by how many clients cover it. This is the one
accuracy cost of dropping alignment, and it is a property of \emph{coverage} rather than of the frame. It is
most visible when the label space is large and the skew extreme, so that each client sees only a sliver of
the classes, and it narrows as each client sees more of the label space, closing well before the data become
uniform. We characterize this gap directly in \cref{sec:experiments}, measured against an alignment-based
reference that closes it only by exposing per-client class summaries in the clear, the exposure the
cryptographic design exists to remove.

\subsection{Local Fine-Tuning}
\label{sec:distill}

Starting from $\thz$, client $j$ fine-tunes its trainable unit on its own data for $\Ksteps$ steps of
ordinary supervised training, producing parameters $\theta_j^{(\Ksteps)}$. The step budget $\Ksteps$ is
deliberately small. This is not an optimisation shortcut but part of the method: a short trajectory keeps
the trainable unit close to $\thz$, and therefore inside the shared frame, so the displacement it produces
is still expressed in the common coordinates. Allowed to train to convergence, each client would drift to
its own local optimum and reintroduce exactly the conflict the shared start was meant to avoid.

At the end of its trajectory the client transmits not its parameters but the displacement
\begin{equation}
\label{eq:displacement}
\disp \;=\; \theta_j^{(\Ksteps)} - \thz ,
\end{equation}
the distance it travelled from the shared start. Because every client departs from the same $\thz$, these
displacements are directly comparable across clients, which is what makes the linear combination that
follows meaningful.

\subsection{Encrypted Linear Aggregation}
\label{sec:aggregation}

The clients hold a joint encryption key produced by distributed key generation, so no single party, and in
particular not the server, ever holds the secret key. Each client encrypts its displacement $\disp$; since
the trainable unit is small relative to the backbone, this is a few tens of ciphertexts. The server fuses the
per-client displacements into a single update of the shared model by the sample-weighted combination
\begin{equation}
\label{eq:aggregate}
\thstar \;=\; \thz + \sum_{j=1}^{\Nc} \wj \, \disp ,
\qquad \wj = \frac{\nj}{\sum_{i} n_i},
\end{equation}
in which each client contributes in proportion to the data $\nj$ behind it. Only the displacements are
encrypted; the weights $\wj$ and the initializer $\thz$ are public and stay in plaintext. Forming
\cref{eq:aggregate} therefore requires only multiplying each ciphertext by the public scalar $\wj$ and
adding the results, with $\thz$ added in plaintext: there is no ciphertext-by-ciphertext multiplication, no
encrypted optimizer state, and no bootstrapping, so the circuit has multiplicative depth one. A threshold of clients
then jointly decrypts $\thstar$ by collective key-switching, and the server never observes a plaintext
update. This is the central design choice: casting aggregation as an operation that is
\emph{linear in the encrypted quantities} yields a server computation that maps onto the least costly operations
a levelled scheme provides, without the encrypted nonlinearities that make encrypted training expensive. The
operations are also fully parallel (the clients fine-tune and encrypt independently, and the
server's sum is an associative reduction), so the end-to-end wall-clock is one client's local computation
plus a single aggregation, independent of $\Nc$. We instantiate \cref{eq:aggregate} in a real multiparty
CKKS implementation and report its measured correctness and cost in \cref{sec:fhe-cost}.

We state precisely why this works, since \cref{eq:aggregate} can be misread. Since
the weights sum to one, the displacement form is algebraically a weighted average of the clients' fine-tuned
units, $\thstar = \sum_j \wj\, \theta_j^{(\Ksteps)}$, which is exactly task-vector merging~\cite{ilharco2023editing}.
This linearity is what keeps the operation inexpensive to encrypt, and it is exact rather than approximate
only because the adapter's down-projection is frozen: were both adapter factors trained, the per-client update
$B_jA_j$ would be bilinear, the average of the factors would not equal the average of the updates, and the
quantity the server can compute at depth one would no longer be the one the method requires. The same average over
units trained independently, from different initialisations, to convergence is the
naive combination that fails under heterogeneity. The average lies in a usable region here because of the protocol around
it: every $\theta_j^{(\Ksteps)}$ starts from the shared $\thz$ and moves only a bounded distance, so all of
them remain in the common frame and their weighted average is itself a reasonable model. The contribution is not
a new aggregation rule but a one-shot protocol under which a linear, and therefore inexpensively encryptable,
aggregation produces a strong shared model. One distinction matters throughout: this depth-one linear family
\emph{is} task-vector merging, and the plain sample-weighted sum ($\lambda{=}1$) is a single point in it ---
the ``naive'' aggregate that heterogeneity destabilises. The client-side selection of \cref{sec:candidates}
does not swap in a different rule; it picks the point in the same family (a softer $\lambda$, or a coverage-
or curvature-reweighted merge) that generalises best, so the linearity, and with it the depth-one encryption,
holds for every candidate.

\subsection{Multi-Candidate Release and Client-Side Selection}
\label{sec:candidates}

A server computing under encryption is blind: it cannot inspect a value, so it cannot adapt the aggregation
rule to the data the way a plaintext aggregator could. We recover that adaptivity without giving the server
the ability to read its inputs, by moving the choice to the clients after decryption. The observation we use
is that several useful aggregation rules are each, on their own, a depth-one function of the encrypted
displacements, so the server evaluates all of them in one pass and the clients select among the results. We
form a small set of candidates, in three tiers:
\begin{itemize}
\item \emph{Reweighted merges} $\thz + (\sum_j \wj g_j\!\odot\!\disp)\oslash(\sum_j \wj g_j)$, where $g_j$ is
  a per-coordinate weight each client forms from its own data: a diagonal Fisher estimate, or --- the variant
  that carries most of our accuracy under skew --- a per-class example count applied to the head rows, so that
  the clients which actually hold a class decide its classifier row and rows no client covers stay at $\thz$.
  This is the merge the client vote most often selects (\cref{sec:superior}). Each client forms the products
  $g_j\!\odot\!\disp$ and $g_j$ \emph{before} encryption; the server adds the two ciphertext stacks; and the
  clients divide the decrypted numerator by the decrypted denominator in plaintext, so the nonlinear division
  never occurs under encryption and the server operation stays a depth-one sum.
\item a \emph{scaling family} $\thz + \lambda\sum_j\wj\disp$ for a public grid of $\lambda$ --- the plain
  weighted aggregate ($\lambda{=}1$) and softer steps toward it --- which trades the shared start against the
  aggregated move and costs nothing extra, since all candidates are affine in the same encrypted sum. This is
  the conservative baseline tier: it discloses nothing beyond the released model (below), and when the data are
  near-uniform it already suffices; under skew it is the reweighted merges that lift the model above it.
\item \emph{leave-one-out} aggregates $\thz + \sum_{i\neq k}\tilde w_i\Delta_i$, one per excluded client $k$,
  with public renormalised weights $\tilde w_i$, used for robustness (\cref{sec:robust}).
\end{itemize}
The server emits the encrypted candidates; a threshold of clients jointly decrypts each one; every client
scores the decrypted candidates on a small held-out slice of its own data; and the federation releases the
candidate with the highest sample-weighted mean held-out accuracy across clients. Selection is therefore
post-decryption and entirely client-side: it adds one threshold decryption per candidate and no homomorphic
multiplications beyond the depth-one aggregate. It could not usefully move under encryption --- it is
data-dependent, scored on held-out client data the server does not have, and homomorphic comparison would
need deep polynomial approximations of the sign function, exactly the circuit cost the design avoids, for a
step the clients perform for free after decryption.

\paragraph{What each candidate reveals.} The candidates differ in what they disclose to the decrypting
clients, and we account for it plainly. The scaling family is collinear,
$\thstar(\lambda)=(1-\lambda)\thz+\lambda\thstar(1)$, so releasing the whole grid reveals nothing beyond the
released model itself. A reweighted merge decrypts its numerator and denominator separately, so it
additionally reveals two \emph{aggregate} sums --- for the count-weighted head, the federation-wide per-class
example totals $\sum_j n_{j,c}$ and the count-weighted head displacement. This is more than the scaling
family discloses, but it is consistent with the threat model of \cref{sec:threat}: both are sums over all
clients, they isolate no individual client, and the server, which never decrypts, learns nothing from them.
A leave-one-out pair, by contrast, does isolate a single client's contribution, and is therefore admitted
only when its robustness use (\cref{sec:robust}) is wanted. Two knobs follow. A deployment content to disclose
nothing beyond the released model restricts the set to the scaling family, forgoing the accuracy the
reweighted merges would have bought; a deployment that wants the merges but not even the aggregate counts
computes the head ratio through a lightweight post-decryption secure division --- one secure reciprocal per
class and a rescale of each head row, negligible beside the homomorphic aggregation --- so that only the
final head rows, and neither sum, are ever revealed. Both prices are small, and neither touches the server's
depth-one operation.

\subsection{Releasing versus Serving the Aggregate}
\label{sec:serve}

The protocol so far produces a single object, the selected aggregate $\thstar$, held under encryption. One
step remains --- making it usable --- and it admits two settings that trade disclosure against cost. The
results of this paper concern the first; the second is available where a deployment may not expose the model
at all.

\emph{Release.} A threshold of clients jointly decrypts $\thstar$, and every party holds the shared model in
the clear. This is the natural output of a federation whose purpose is a common model, and it is the setting
the accuracy, robustness, and membership-inference results below assume. Its only residual surface is the
released model itself (\cref{sec:security,sec:mia}).

\emph{Serve.} The aggregate is never decrypted. It stays in ciphertext and answers queries under encryption,
so that no party --- client or server --- ever holds the model in the clear, and a client learns only the
labels it requests. This fits deployments in which the fine-tuned model is itself a regulated or proprietary
asset that may not be released even to its builders. It is affordable for the same reason the aggregation is:
the trainable unit is \emph{linear in the frozen features}. For a query $x$ the public backbone runs in
plaintext to features $\featx(x)$; the encrypted head maps them to encrypted logits by the very operation the
aggregation uses --- a plaintext vector against ciphertext weights, at multiplicative depth one --- so no
nonlinearity ever meets the encrypted head. The predicted label is the $\arg\max$ of those logits, and a
threshold of clients decrypts that label alone.

\paragraph{Why the label and not the logits.} Returning logits would defeat the purpose. The head is a linear
map on features the client already holds, so a client that reads logits for feature vectors of its choosing
recovers the head exactly, by solving a linear system in as many queries as the feature
dimension~\cite{tramer2016stealing}. Releasing only the $\arg\max$ denies this: it reduces an adversary to
label-only boundary search~\cite{lowdmeek2005,choquettechoo2021labelonly}, which recovers decision boundaries
but never the head's coordinates, and at far greater query cost. Where honest use is intrinsically
high-volume, so that no per-client budget separates it from extraction, the returned labels can additionally
carry calibrated noise --- report-noisy-max~\cite{mcsherry2007mechanism}, folded into the smudging the
threshold decryption already performs so that no sub-threshold coalition can remove it --- yielding a
differential-privacy bound on cumulative leakage. This is a query-budget bound, not a cryptographic one: the
aggregate remains recoverable in principle, only expensively.

\paragraph{Cost, and who serves.} Serving carries one price the rest of the protocol avoids: the encrypted
$\arg\max$ is a homomorphic comparison circuit~\cite{cheon2020comparison}, the one operation in the method
that is not depth-one, and it is paid per query rather than once. Evaluated naively, as a sequential fold of
$\Cc-1$ pairwise comparisons, it costs on the order of minutes for the larger label spaces --- roughly twenty
minutes at $\Cc{=}77$ in our CPU implementation. This is a loose upper bound, not a floor: the $\arg\max$ is
ordinary homomorphic computation, and the standard optimisation applies. Packing the $\Cc$ logits into one
ciphertext and reducing them with a log-depth SIMD tournament cuts the sequential comparisons from $\Cc-1$ to
$\lceil\log_2\Cc\rceil$ --- a fourteen-fold reduction at $\Cc{=}100$ --- and secure-transformer systems built
on exactly this reduction report a full-vocabulary $\arg\max$ in well under a second on a
GPU~\cite{zhang2025nexus}. None of this is specific to the threshold setting: homomorphic evaluation under a
collectively generated key is identical to the single-key case, so these results carry over directly, the
multiparty scheme differing only in its distributed key generation and its threshold decryption of the final
label. The per-query cost is therefore minutes at worst and sub-second at best, which is why Serve suits
low-throughput or governance deployments and Release the rest. The party that answers queries is not thereby
trusted with anything: holding only the public key, the evaluation keys, and the ciphertext, it computes but
cannot decrypt, and so learns neither the model nor any client's data. It is trusted only for availability
and, where honest computation must be certified, through verifiable homomorphic
encryption~\cite{viand2023verifiable}; the decryption threshold is unchanged, since the aggregate becomes
plaintext only to a quorum, exactly as under Release.

\subsection{Threat Model}
\label{sec:threat}

The participants are $\Nc$ clients and a single aggregation server. The server is honest-but-curious: it
follows the protocol but may try to infer private information from anything it observes. During aggregation
it sees only ciphertexts under the collective key, and as established in \cref{sec:mhe} it cannot decrypt
them. A minority of clients may likewise be honest-but-curious and may collude, but any coalition below the
decryption threshold learns nothing about another client's displacement, since recovering a plaintext
requires the cooperation of the threshold. This is a stronger position than secure-aggregation protocols,
where the revealed sum has itself been shown to leak across rounds~\cite{so2023securing,kerkouche2023client},
because in our one-shot protocol the server never obtains a plaintext sum at all.

We are explicit about one boundary of this model, which is the training-time/inference-time split of
\cref{sec:related-leakage}. The protocol removes the training-time surface entirely: no gradient, update,
or data-derived summary is ever observable in plaintext, in any round, because there is only one round and
its single artifact is encrypted. Every participant does obtain the released model $\thstar$ in the clear
after threshold decryption, by design, since the purpose of the protocol is to hand the parties a shared
model. Any inference a participant then performs on $\thstar$ against another participant's data is
therefore not an attack the protocol can prevent but a property of releasing a useful shared model at all,
and we measure it directly in \cref{sec:mia} under exactly this fellow-participant adversary. Where a
deployment may not expose even the released model, the Serve setting of \cref{sec:serve} withholds $\thstar$
altogether, leaving a participant only the labels of its own queries, at the per-query serving cost discussed
there; the contribution-level guarantee stated next is identical in both settings. The
cryptographic guarantee concerns the \emph{contributions}: the server and any sub-threshold coalition observe
only ciphertexts. This is the sense in which the protection target differs from differentially private
one-shot methods, which perturb the released artifact to protect it against all parties, including
participants, at a cost to its accuracy. Robustness to actively
malicious or Byzantine clients is outside the cryptographic guarantee; the multi-candidate mechanism of
\cref{sec:candidates} provides a lightweight measured defense (\cref{sec:robust}), and stronger guarantees we
discuss as future work in \cref{sec:scope}.

\subsection{Security}
\label{sec:security}

The model updates are protected cryptographically: the aggregation of \cref{eq:aggregate} runs entirely
under multiparty CKKS with threshold decryption, so the displacements $\disp$ are never exposed to the
server or to any minority of clients, and they are protected without any perturbation. The cryptographic
layer adds no accuracy cost. This is the central privacy claim, and the point on which the method departs
from differentially private one-shot federated learning, which injects noise into the updates it ultimately
releases and so pays a utility tax on the very thing it ships.

We make this precise. Let the view of the honest-but-curious server be everything it receives: the public
collective key, the public weights $\{\wj\}$ and initializer $\thz$, the client ciphertexts
$\{\mathsf{Enc}(\disp)\}_{j=1}^{\Nc}$, and the homomorphically computed $\mathsf{Enc}(\thstar)$. We write
$\mathsf{view}_{\mathrm{srv}}(\{\Dj\})$ for this view as a function of the private datasets.

\begin{proposition}
\label{prop:server}
Under the IND-CPA security of the multiparty CKKS scheme and its $t$-out-of-$\Nc$ threshold decryption, for
any honest-but-curious server colluding with up to $t-1$ clients there is a probabilistic polynomial-time
simulator $\mathcal{S}$, given only the public parameters, such that
$\mathsf{view}_{\mathrm{srv}}(\{\Dj\}) \stackrel{c}{\approx} \mathcal{S}$. In particular the server's view is
computationally independent of the private datasets $\{\Dj\}$ and the displacements $\{\disp\}$, and the
server cannot recover any plaintext, including $\thstar$, without a decryption quorum.
\end{proposition}

\begin{proof}[Proof sketch]
The simulator outputs the public parameters together with $\Nc$ encryptions of zero under the collective
key. By IND-CPA each genuine $\mathsf{Enc}(\disp)$ is computationally indistinguishable from
$\mathsf{Enc}(0)$; a hybrid over the $\Nc$ ciphertexts replaces them one at a time, so the whole collection
is indistinguishable from the simulated one and carries no information about $\{\disp\}$ or the data they
derive from. The aggregate $\mathsf{Enc}(\thstar)$ is a deterministic function of those ciphertexts and the
public $\{\wj\},\thz$, so it adds nothing. Decryption needs $t$ shares of a secret key that is never
reconstructed in one place; the server with at most $t-1$ clients holds too few shares to decrypt.
\end{proof}

\Cref{prop:server} pins down where information can flow. Because there is no alignment phase, no data-derived
quantity leaves a client before encryption, so the protocol exposes private information through exactly one
channel, intrinsic to the task rather than an artifact of our construction: the released model itself. After
threshold decryption every client holds $\thstar$ in the clear, and white-box access to the jointly built
model is unavoidable in any protocol whose purpose is to hand the parties a shared model. The residual
information that model carries about training data is what we quantify with membership inference in
\cref{sec:mia}.
